\section*{Chapter \ref{ch:g11:money-and-banking} --- Money and Banking}
% figdata: figdata/grade-11/13-money-and-banking/out/usa_money.csv; Brenland's banks invented.

\begin{solution}{exo:g11:money-and-banking:1}
Assets: \omterm{def:g11:money-and-banking:reserves}{reserves} (cash) $+500$. Liabilities: deposits $+500$. \omterm{def:g11:money-and-banking:broad-money}{Broad money} is unchanged: the household holds 500 less in notes and 500 more in deposits.
\end{solution}

\begin{solution}{exo:g11:money-and-banking:2}
Assets: loans $+1000$. Liabilities: deposits $+1000$. \omterm{def:g11:money-and-banking:broad-money}{Broad money} rises by 1000.
\end{solution}

\begin{solution}{exo:g11:money-and-banking:3}
Assets: bonds $+200$. Liabilities: deposits $+200$. \omterm{def:g11:money-and-banking:broad-money}{Broad money} rises by 200: buying an asset with a new deposit creates money just as a loan does.
\end{solution}

\begin{solution}{exo:g11:money-and-banking:4}
Narrow money: $3426/847 \approx 4.0$. \omterm{def:g11:money-and-banking:broad-money}{Broad money}: $14\,836/7294 \approx 2.0$.
\end{solution}

\begin{solution}{exo:g11:money-and-banking:5}
$40/500 = 8\,\%$; deposits $500 - 40 = 460$.
\end{solution}

\begin{solution}{exo:g11:money-and-banking:6}
$(1 + 0.2)/(0.2 + 0.05) = 4.8$; \omterm{def:g11:money-and-banking:broad-money}{broad money} $4.8 \times 100 = 480$.
\end{solution}

\begin{solution}{exo:g11:money-and-banking:7}
$c = 0$: $1/0.1 = 10$. $c = 0.1$: $1.1/0.2 = 5.5$. Cash taken out of banks is not lent again, so each round of lending creates fewer new deposits.
\end{solution}

\begin{solution}{exo:g11:money-and-banking:8}
Assets: loans $-50$. Liabilities: deposits $-50$. \omterm{def:g11:money-and-banking:broad-money}{Broad money} falls by 50: repaying a loan destroys money.
\end{solution}

\begin{solution}{exo:g11:money-and-banking:9}
If banks lend all excess \omterm{def:g11:money-and-banking:reserves}{reserves}: the multiplier is $1.1/0.2 = 5.5$, and \omterm{def:g11:money-and-banking:broad-money}{broad money} rises by 550. If they hold them: bonds are swapped for \omterm{def:g11:money-and-banking:reserves}{reserves} inside the banks,
no one's deposit changes, and \omterm{def:g11:money-and-banking:broad-money}{broad money} does not change.
\end{solution}

\begin{solution}{exo:g11:money-and-banking:10}
Before: $60/1000 = 6\,\%$. After: capital $35$, assets $975$, ratio $35/975 \approx 3.6\,\%$.
\end{solution}

\begin{solution}{exo:g11:money-and-banking:11}
7.8 in 2000, 8.6 in 2007, 4.7 in 2008, 2.9 in 2014. It fell most in 2008, from 8.6 to 4.7.
\end{solution}

\begin{solution}{exo:g11:money-and-banking:12}
The model predicts a stable ratio of broad to narrow money; the chart shows the ratio near 8 until 2007, then halving in one year and falling to 2.9 by 2014, while narrow
money quadrupled. The assumption that failed is that banks lend every reserve beyond what they must keep: they held the new \omterm{def:g11:money-and-banking:reserves}{reserves}, which earned interest at the
central bank, while creditworthy borrowers were scarce.
\end{solution}

\begin{solution}{exo:g11:money-and-banking:13}
If the other stays, staying gives 110 and running 100: stay. If the other runs, staying gives 50 and running 75: run. So (stay, stay) and (run, run) are both equilibria.
With insurance paying 100 to a depositor who stays, staying gives 100 against 75 when the other runs, and 110 against 100 when the other stays: staying is best whatever
the other does, and the run equilibrium disappears.
\end{solution}

\begin{solution}{exo:g11:money-and-banking:14}
The bank owes $90 \times 1.02 = 91.8$. Safe asset: $105 - 91.8 = 13.2$ for the owners. Risky asset: $130 - 91.8 = 38.2$ or, with assets of 80, nothing (their loss stops
at their capital), an expected $19.1$; the insurer pays $91.8 - 80 = 11.8$ half the time, an expected cost of $5.9$. Both assets are worth 105 on average, but the owners
gain from choosing risk, at the insurer's expense: \omterm{def:g11:money-and-banking:deposit-insurance}{deposit insurance} creates \omterm{def:g11:asymmetric-information:moral-hazard}{moral hazard}, which capital rules try to contain.
\end{solution}

\begin{solution}{exo:g11:money-and-banking:15}
A loss of 40 wipes out capital. A loss $L$ breaches 6\,\% when $(40 - L)/(500 - L) < 0.06$, that is $40 - L < 30 - 0.06L$, or $L > 10/0.94 \approx 10.6$.
\end{solution}

\begin{solution}{pb:g11:money-and-banking:1}
\textbf{1.} Assets $20 + 30 + 170 = 220$; capital $220 - 200 = 20$; ratio $20/220 \approx 9.1\,\%$.

\textbf{2.} Loans $+10$, deposits $+10$. Assets 230, capital 20, ratio $20/230 \approx 8.7\,\%$.

\textbf{3.} Bren Savings Bank: \omterm{def:g11:money-and-banking:reserves}{reserves} $-10$, deposits $-10$. The other bank: \omterm{def:g11:money-and-banking:reserves}{reserves} $+10$, deposits $+10$. \omterm{def:g11:money-and-banking:broad-money}{Broad money} is unchanged by the payment.

\textbf{4.} Loans $-5$, deposits $-5$: \omterm{def:g11:money-and-banking:broad-money}{broad money} falls by 5.

\textbf{5.} The bank now has \omterm{def:g11:money-and-banking:reserves}{reserves} 10, bonds 30, loans 175 and deposits 195. The loss leaves loans of 167, assets of 207 and capital of 12: a ratio of
$12/207 \approx 5.8\,\%$.

\textbf{6.} \omterm{def:g11:money-and-banking:reserves}{Central bank money} is cash plus \omterm{def:g11:money-and-banking:reserves}{reserves}: $0.1D + 0.1D = 200$, so deposits are 1000, cash 100, \omterm{def:g11:money-and-banking:reserves}{reserves} 100, and \omterm{def:g11:money-and-banking:broad-money}{broad money} $100 + 1000 = 1100$.

\textbf{7.} $1100/200 = 5.5$, which is $(1 + 0.1)/(0.1 + 0.1)$.

\textbf{8.} Deposits $400/0.2 = 2000$, cash 200, \omterm{def:g11:money-and-banking:broad-money}{broad money} 2200.

\textbf{9.} The banks swap bonds for \omterm{def:g11:money-and-banking:reserves}{reserves} and lend nothing more: deposits stay at 1000 and cash at 100, so \omterm{def:g11:money-and-banking:broad-money}{broad money} stays at 1100; \omterm{def:g11:money-and-banking:reserves}{reserves} are now 300.

\textbf{10.} $2200/400 = 5.5$ if banks lend; $1100/400 = 2.75$ if they hold.

\textbf{11.} Neither exactly: \omterm{def:g11:money-and-banking:broad-money}{broad money} would have quadrupled with narrow money in the first regime, and not grown at all in the second. It doubled, and the multiplier
halved: much of the new \omterm{def:g11:money-and-banking:reserves}{reserves} was held.

\textbf{12.} \omterm{def:g11:money-and-banking:reserves}{Reserves} held at the central bank earn interest and carry no risk; creditworthy borrowers may be few in a recession; banks short of capital cannot add risky
loans; rules may require them to hold liquid assets.

\textbf{13.} Every new loan adds to its risk-weighted assets but not to its capital, so its \omterm{def:g11:money-and-banking:capital-ratio}{capital ratio} falls. Once the ratio is at its minimum, the bank cannot lend more
without raising capital, whatever its \omterm{def:g11:money-and-banking:reserves}{reserves}.

\textbf{14.} \omterm{def:g11:money-and-banking:reserves}{Reserves} plus bonds, $20 + 30 = 50$: 25\,\% of deposits.

\textbf{15.} It loses 1 of capital for each 1 paid out.

\textbf{16.} Capital of 20 is gone after 20 more is paid out from loan sales, which uses loans recorded at 40. Total withdrawals: $50 + 20 = 70$, or 35\,\% of deposits.

\textbf{17.} The first 70 are paid in full. After that the bank has loans recorded at 130 that fetch only 65 in a quick sale, against deposits of 130: if all the remaining
depositors run, they get half their money. Waiting is costly if the others run, so each depositor has a reason to run first.

\textbf{18.} Insured depositors are paid in full whatever happens, so none gains by running; a run does not start.

\textbf{19.} The bank could raise cash against its loans without selling them at a loss: it pays every depositor who asks and keeps its capital, so running gains nothing.

\textbf{20.} After the central bank doubles its money to 400, \omterm{def:g11:money-and-banking:broad-money}{broad money} is 2200 if banks lend all excess \omterm{def:g11:money-and-banking:reserves}{reserves} and stays at 1100 if they hold them; the Bren Savings
Bank can pay out 35\,\% of its deposits before its capital is gone.
\end{solution}
